% Teaching manuscript. Edit this TeX source; Markdown is compatibility navigation only.

\StudyChapter{CS5222-Tutorial02}{Tutorial 2 · Delay and traceroute｜先定义计时事件，再按顺序相加}
\StudyTarget{CS5222-Tutorial02__tutorial-2-delay-and-traceroute}


\StudyRef{CS5222-Chapter01}{delay}{Chapter 1} · \StudyRepoRef{CS5222/course-notes/README.md}{}{课程目录} · \StudyRef{CS5222-Tutorial01}{root}{Tutorial 1} · \StudyRef{CS5222-Assignment01}{root}{Assignment 1}


依据当前 原题PDF（3页） 与 教师解答PPT（13页），按Q1–5讲解。难点是计时边界与转发模型；优先级均为核心必会，Q1/Q3中等偏难，Q2/Q4中等，Q5需实际观察与解释。优先级依据当前Tutorial。


\textbf{基础复习（可跳过）：} \StudyRef{Foundations-NetworkBasics}{units}{单位}、\StudyRef{Foundations-NetworkBasics}{timeline}{首位与末位时间轴}。所有计算约定Mbps=10\textsuperscript{6}bit/s、kbps=10\textsuperscript{3}bit/s。






\StudyHeading{1}{Q1. VoIP｜语音先攒成一包，还要算这段等待}
\StudyTarget{CS5222-Tutorial02__q1}
\StudyTarget{CS5222-Tutorial02__q1-voip}


\StudyImage{CS5222/course-notes/assets/tutorial02-01.png}


\textbf{完整条件 / Conditions:} A以128kbps实时生成语音比特，每64byte成一包，收集完成立即发送；链路4Mbps、传播8ms，B收完整包后还原模拟语音。问生成一个bit到它被解码的时延。 / A generates voice at 128 kbps, groups it into 64-byte packets and immediately transmits each completed packet over a 4-Mbps link with 8-ms propagation. B decodes after receiving the complete packet. Find generation-to-decoding delay.


先按包首位及连续流体近似计算：包长512bit；收集一包 \(512/128000=4\)ms；把包送上链路 \(512/4000000=0.128\)ms；传播8ms。首位从生成到可开始播放的延迟为 \textbf{12.128ms，选D}，与Tutorial 2教师解答，第5张幻灯片结论一致。


这里“任意一位”还要区分读取与播放这两个时刻。一个包内越晚产生的bit等待成包越短；若问“该bit到达整包缓冲、可被读取”的延迟，数值随位次变化，约从12.128ms降到8.128ms。若接收端在包到齐后，\textbf{按与源相同128kbps的节奏顺序恢复/播放语音}，后面的bit在接收端也会多等相同位次时间，因此每个bit的生成到播放时延相同，都是12.128ms。教师“all bits”的说法应连着这个播放约定读。


设第j位相对首位晚产生 \(\tau_j\)，包到达时刻T=12.128ms；可读取延迟为T\ensuremath{-}\ensuremath{\tau}\textsubscript{j}，顺序播放时刻为T+\ensuremath{\tau}\textsubscript{j}，播放延迟则(T+\ensuremath{\tau}\textsubscript{j})\ensuremath{-}\ensuremath{\tau}\textsubscript{j}=T。因此，答题前要确定问的是可读取时刻还是播放时刻。


\textbf{English answer:} Under the instructor's packetization and source-rate sequential playout model, the delay is 4+0.128+8=12.128 ms for every bit. Complete-packet availability alone gives a position-dependent delay.


\textbf{变式 / Transfer:} 包改128byte，其他条件不变，首位/同速播放时延？ / Double packet size to 128 bytes, keeping all rates and propagation unchanged.



\begin{StudySolution}{答案 / Answer}


8+0.256+8=16.256ms。 / 16.256 ms. 较大包减少每份数据的头部比例，但收集等待变长；本题没有给头长，不计算额外百分比。


\end{StudySolution}






\StudyHeading{1}{Q2. Three store-and-forward links｜每一跳的发送与传播都要数}
\StudyTarget{CS5222-Tutorial02__q2}
\StudyTarget{CS5222-Tutorial02__q2-three-store-and-forward-links}


\textbf{条件 / Conditions:} 长L bit的包经三条链路，两台存储转发交换机；第i链路长d\textsubscript{i} m、传播速度s\textsubscript{i} m/s、发送速率R\textsubscript{i} bit/s，每台交换机处理时延dproc，无排队。 / An L-bit packet crosses three store-and-forward links and two switches. Link i has length d\textsubscript{i}, propagation speed s\textsubscript{i} and rate R\textsubscript{i}; each switch adds dproc and there is no queueing.


a 把事件从左到右排：链路1发送/传播 \ensuremath{\to} 交换机1处理 \ensuremath{\to} 链路2发送/传播 \ensuremath{\to} 交换机2处理 \ensuremath{\to} 链路3发送/传播。



\[
d_{end}=\sum_{i=1}^3\frac L{R_i}+\sum_{i=1}^3\frac{d_i}{s_i}+2d_{proc}\quad\text{（选B）}.\tag{T2.1}
\]


b 给定L=1500byte=12000bit，各R=2Mbps，各s=2.5\ensuremath{\times}10\textsuperscript{8}m/s，距离5000/4000/1000km，每交换机3ms：


\begingroup
\small\setlength{\tabcolsep}{3pt}
\renewcommand{\arraystretch}{1.22}
\begin{longtable}{@{}>{\raggedright\arraybackslash}p{\dimexpr0.299276\linewidth-1.33333\tabcolsep\relax}>{\raggedright\arraybackslash}p{\dimexpr0.382200\linewidth-1.33333\tabcolsep\relax}>{\raggedright\arraybackslash}p{\dimexpr0.318524\linewidth-1.33333\tabcolsep\relax}@{}}
\toprule
部分 & 计算 & 结果 \\
\midrule\endfirsthead
部分 & 计算 & 结果 \\
\midrule\endhead
\bottomrule\endfoot
三段发送 & 每段12000/(2\ensuremath{\times}10\textsuperscript{6}) & 6+6+6ms \\
三段传播 & 每段先km\ensuremath{\to}m再除s & 20+16+4ms \\
两次处理 & 2\ensuremath{\times}3ms & 6ms \\
总计 & 18+40+6 & \textbf{64ms，选B} \\
\end{longtable}
\endgroup


\textbf{English answer:} Sum the three serialization times, three propagation times and two processing times. The numerical total is 64 ms.


这里包长L在公式中必须是bit；若直接使用1500需在分子明确乘8。处理是两台交换机，不是三条链路所以乘3。






\StudyHeading{1}{Q3. Immediate bit forwarding｜不用每个中间节点再等整包}
\StudyTarget{CS5222-Tutorial02__q3}
\StudyTarget{CS5222-Tutorial02__q3-immediate-bit-forwarding}


\StudyImage{CS5222/course-notes/assets/tutorial02-02.png}


\textbf{条件 / Conditions:} 沿用Q2，R\textsubscript{1}=R\textsubscript{2}=R\textsubscript{3}=R、dproc=0，无排队，中间节点每收到一bit立即转发，不等整包。 / Use the Q2 path, equal rates and zero processing; each bit is forwarded immediately with no whole-packet waiting or queueing.


首位经过各传播段，末位相对首位只落后源端一次L/R。故总为 \(L/R+\sum_i d_i/s_i\)。带Q2数字为6+20+16+4=\textbf{46ms，选D}，Tutorial 2教师解答，第8张幻灯片相同。


这不表示物理链路“传输时间消失”，而是各段的发送过程可以流水重叠，末位到达总时间只额外保留一次包序列长度。不同速率、实际cut-through启动/头部检查有其他限制；不要把这个理想模型直接当所有交换机行为。


\textbf{English answer:} Equal-rate immediate bit forwarding gives L/R plus all propagation delays, hence 46 ms. Serialization on successive links overlaps.


\textbf{变式 / Transfer:} 四条相同速率链路，每条传播2ms，源包发送3ms，即时转发和整包存储转发分别多久？ / Four equal-rate links each propagate for 2 ms; serialization takes 3 ms. Compare immediate bit forwarding with store-and-forward for one packet.



\begin{StudySolution}{答案 / Answer}


即时3+4\ensuremath{\times}2=11ms；存储转发4\ensuremath{\times}3+4\ensuremath{\times}2=20ms，忽略排队/处理。 / 11 ms versus 20 ms, with no queueing or processing.


\end{StudySolution}






\StudyHeading{1}{Q4. Queueing｜只算轮到你之前的等待}
\StudyTarget{CS5222-Tutorial02__q4}
\StudyTarget{CS5222-Tutorial02__q4-queueing}


\textbf{条件 / Conditions:} 新包到达时，出口上一个1500byte包已发一半，还有4包等待，均1500byte，FIFO，R=2Mbps。求新包排队时延，再推广为包长L bit、当前包已发送x bit、另有n包排队。 / A new packet finds one equal packet half transmitted and four equal packets queued, FIFO, L=1500 bytes and R=2 Mbps. Find queueing delay, then generalize to x bits already transmitted and n waiting packets.


前面剩4.5包，\(4.5\times12000/(2\times10^6)=0.027\)s=\textbf{27ms，选A}。一般式为



\[
d_{queue}=\frac{(L-x)+nL}R.\tag{T2.2}
\]


\textbf{English answer:} Queueing delay is 27 ms, or ((L\ensuremath{-}x)+nL)/R. It excludes transmission of the arriving packet itself.


若问它\textbf{完成发送}的时间，则另加L/R；若问到目的地还要加传播及后续节点条件。若出口本来空闲，就不是这条“有一个正在发送包”的场景，等待为0。


\textbf{变式 / Transfer:} 同速同包长，当前包已发75\%，队列前面还有2包，新包排队多久？ / The current packet is 75\% transmitted and two packets are queued ahead. Find the new packet's queueing delay.



\begin{StudySolution}{答案 / Answer}


2.25\ensuremath{\times}6=13.5ms。 / 13.5 ms. 不是加上新包自己的6ms变成19.5ms。


\end{StudySolution}






\StudyHeading{1}{Q5. Traceroute｜把真实测量和推断分开写}
\StudyTarget{CS5222-Tutorial02__q5}
\StudyTarget{CS5222-Tutorial02__q5-traceroute}


\textbf{原任务 / Task:} 对cityu.edu.hk运行traceroute，解释每列、星号和后跳RTT较小的原因。 / Run traceroute to cityu.edu.hk and explain its columns, asterisks, and non-monotonic RTTs.


macOS可运行原命令 \StudyInline{traceroute cityu.\allowbreak{}edu.\allowbreak{}hk}。下方2026-09-23的观测使用：



\begin{StudyCode}
traceroute -n -m 12 -q 3 -w 1 cityu.edu.hk
\end{StudyCode}


\StudyInline{-n}保留数值IP，\StudyInline{-m 12}最多12跳，\StudyInline{-q 3}每跳3次探测，\StudyInline{-w 1}每次最多等1秒。这些参数用于限制本次探测的长度与等待时间，并非题目指定值。输出与时间见2026-09-23观察，采集脚本为observe\_network.py。


每行第一项为TTL/hop编号，随后为响应地址和三次RTT（ms）；某次无及时响应显示\StudyInline{*}。地址识别的是响应接口，复杂网络里不保证完整暴露每一台中间设备。macOS常用UDP探测和ICMP响应；Windows \StudyInline{tracert}的实现细节不同，不应机械照搬命令选项。


星号可以来自过滤、应答限速、探测/响应丢失、超时等。Tutorial 2教师解答，第12张幻灯片强调防火墙是常见原因，但仅凭星号不能确定具体根因。后跳RTT小于前跳，可能因为不同时刻队列变化、回程路由不同、路由器处理探测优先级不同；RTT也不是沿途正向传播的逐段累加测量。


\textbf{English answer:} Columns identify probe hop, responding address and round-trip times. An asterisk means no response within the timeout. Different probes and return paths can produce non-monotonic RTTs.


\textbf{已保存的实际观察（2026-09-23 01:19，UTC+8）：} 域名解析有两个地址，traceroute选用45.60.199.218；第1跳三个RTT约4.352/3.957/4.301ms，第7跳有星号，第11跳第三次升到192.365ms。在最多12跳内没有看到目的地址响应，故这不是完整路径。进程正常结束也不表示探测已到终点。同一跳有不同地址响应，不足以断言一条固定唯一线路。


你自己的观测应另记时间、探测参数和实际输出；换时间或网络，路径和RTT可能变化。



\StudyHeading{1}{回查与巩固}
\StudyTarget{CS5222-Tutorial02___1}


原题p1=Q1/Q2，p2=Q3/Q4，p3=Q5。Tutorial 2教师解答，第1–4张幻灯片为四类时延/车队背景（见Chapter1），5=Q1，6–7=Q2，8=Q3，9=Q4，10–11=traceroute课堂示例，12–13=Q5解释。全部已对应，原教师结论与本文补充条件分开。


计算复核：\StudyRepoRef{CS5222/course-notes/network-example-results.json}{}{network-example-results.json}。复习可用\href{https://crazyshout.github.io/micro-course/cards.html\#CS5222-N025}{现有时延卡N025}与\href{https://crazyshout.github.io/micro-course/cards.html\#CS5222-N080}{原题卡N080}，根据本题错误位置，复习单位、时序、转发或排队的相关概念。
